% =============================================================================
%  C1: Closure-operator unification of viability — skeleton paper
%  STATUS: UNPROVEN. Two cases partially proven in literature; MCC novel.
% =============================================================================
\documentclass[11pt]{article}

\usepackage[utf8]{inputenc}
\usepackage[T1]{fontenc}
\usepackage[expansion=false]{microtype}
\usepackage{amsmath, amssymb, amsthm, mathtools}
\usepackage[margin=1in]{geometry}
\usepackage{enumitem}
\usepackage[hidelinks]{hyperref}

\theoremstyle{plain}
\newtheorem{theorem}{Theorem}
\newtheorem*{namedconjecture}{Conjecture}
\theoremstyle{definition}
\newtheorem{definition}[theorem]{Definition}

\title{Closure-operator Unification of the Four Viability Formalisms \\
       \large (UNPROVEN --- skeleton paper for C1)}
\author{Liam McGuire\thanks{Unaffiliated researcher, Seattle, USA. liammcg44@gmail.com}}
\date{\today}

\begin{document}
\maketitle

\begin{abstract}
This paper is a \textbf{skeleton} that records C1 from the scaffolding
paper~\cite{emergent_systems_2026}. \textbf{No proof is attempted.}
Two of the four cases (autopoietic, RAF) are mostly proven in the
literature; the Markov-blanket case is genuinely open; the MCC case
requires a novel reformulation. The result is UNPROVEN.
\end{abstract}

\section{Statement (UNPROVEN)}

\begin{quote}
\emph{``C1 (closure-operator unification of viability): each of the
four viability formalisms --- Markov-blanket viability, autopoietic
closure, RAF set membership, and MCC --- corresponds to a closure
operator $c_\bullet$ on a partial order of supporting sets, in the
sense that an entity passes the filter iff its supporting set is
closed under $c_\bullet$.''}
\end{quote}
\noindent (Scaffolding paper~\cite{emergent_systems_2026},
\S\texttt{sec:c1}.)

\section{Plan of attack (no proof attempted)}

Four sub-cases, in order of difficulty.

\begin{enumerate}[leftmargin=2em,itemsep=4pt]
  \item \textbf{RAF case (mostly proven in literature).} The RAF
    property is exactly closure under reflexive autocatalysis and
    food-generation. Hordijk \& Steel's polynomial-time
    algorithm~\cite{hordijk2004raf} is a Kleene-iteration to the
    greatest fixed point. The closure operator is
    $c_{\mathrm{RAF}}(R) = R \cup \{r : r \text{ catalysed and reactants in } R\}$.
    Citing this is the work.
  \item \textbf{Autopoietic case (mostly proven in literature).}
    Beer's formalisation of Game-of-Life
    gliders~\cite{beer2014glider,beer2015autopoiesis} uses preimage-
    closure on the reaction network. Translating to the closure-
    operator vocabulary is the work.
  \item \textbf{Markov-blanket case (UNPROVEN; genuinely open).}
    Sparse-coupling implies blanket structure only under additional
    conditions. The closure operator (if it exists) is not obvious.
    Heins \& Da Costa's conditions are the closest existing
    formalisation.
  \item \textbf{MCC case (UNPROVEN; novel work required).} The minimal
    criterion is a predicate, not obviously a closure. The proposed
    reformulation is ``smallest set of solver-and-problem pairs
    satisfying mutual qualification.'' Verifying this is a genuine
    closure operator is the load-bearing novel work.
\end{enumerate}

\section{Status and follow-up}

\textbf{Current status: UNPROVEN.} Two cases citable, two cases open.
A proof of C1 simplifies the construction of the $\mathrm{Decomp}$
functor in \texttt{n1\_necessity/}.

\bibliographystyle{plain}
\bibliography{references}

\end{document}
